% ---------------------------------------------------------------------
%  A certified multilevel lower bound for the Galerkin prediction error
%  of a log-singular Toeplitz form with a prime-power comb.
%
%  Standalone draft.  Self-contained preamble, no external \input:
%  the paper is deliberately independent of any surrounding document
%  set and uses no vocabulary beyond standard numerical analysis.
% ---------------------------------------------------------------------
\documentclass[11pt,a4paper]{article}
\usepackage[T1]{fontenc}
\usepackage[utf8]{inputenc}
\usepackage[english]{babel}
\usepackage{lmodern}
\usepackage[a4paper,margin=2.6cm]{geometry}
\usepackage{amsmath,amssymb,amsthm}
\usepackage{mathtools}
\usepackage{booktabs,array,tabularx}
\usepackage{enumitem}
\usepackage{microtype}
\usepackage[hidelinks]{hyperref}
\emergencystretch=3em

\theoremstyle{plain}
\newtheorem{theorem}{Theorem}[section]
\newtheorem{proposition}[theorem]{Proposition}
\newtheorem{lemma}[theorem]{Lemma}
\newtheorem{corollary}[theorem]{Corollary}
\theoremstyle{definition}
\newtheorem{definition}[theorem]{Definition}
\newtheorem{convention}[theorem]{Convention}
\newtheorem{remark}[theorem]{Remark}

\newcommand{\R}{\mathbb{R}}
\newcommand{\Z}{\mathbb{Z}}
\newcommand{\N}{\mathbb{N}}
\newcommand{\T}{\mathsf{T}}
\newcommand{\Tm}[1]{T_{#1}}
\newcommand{\lmin}{\lambda_{\min}}
\newcommand{\lmax}{\lambda_{\max}}
\newcommand{\ran}{\operatorname{ran}}
\newcommand{\cond}{\operatorname{cond}}
\newcommand{\shat}{\widehat{s}}
\newcommand{\PSD}{\succeq 0}
\newcommand{\PD}{\succ 0}
\DeclareMathOperator{\diag}{diag}
\setlist{itemsep=2pt,topsep=3pt}
\newcolumntype{Y}{>{\raggedright\arraybackslash}X}

\title{A certified multilevel lower bound for the Galerkin\\
prediction error of a log-singular Toeplitz form\\
with a prime-power comb}
\author{Stefan Hamann}
\date{\today}

\begin{document}
\maketitle

\begin{abstract}
\noindent
For a reflection-odd compression $A$ of a finite Toeplitz form and a data vector $b$,
the prediction defect $\varepsilon=1-b^{\T}A^{-1}b$ is the Galerkin error of a
piecewise-constant approximation to a fixed function, and bounding it from
\emph{below} is the hard direction. Our result is a single inequality: on a nested
ladder of resolutions the increment of the saturation telescope admits a certified
lower bound by one explicit test vector, in closed form. The mechanism is a change of
direction --- the increment is a residual measured in the \emph{inverse} energy norm,
hence a \emph{maximum} over test vectors rather than a minimum, so one test vector
already suffices and the coarse-level block enters only through its sign, never
through its smallest eigenvalue, its condition number, or a strengthened
Cauchy--Schwarz constant. The proof combines Schur complement monotonicity, Parseval
against a certified majorant, and L\"owner's antitony of the inverse; the nesting the
ladder needs is proved as an identity for any lag family generated by a fixed kernel.
Two features of the statement are recorded up front rather than in a remark. First, a
scope: the numerator of the certified rung is evaluated at the coarse Galerkin
solution carried along the ladder, which is the information an a posteriori estimator
has, so the bound is a bound \emph{given} the ladder. Second, a measured negative
result: the solution-free variant of the same rung fails, not by a small margin but by
cancellation --- the coupling term annihilates the explicit high-pass part of the data
instead of perturbing it, and a norm-only Cauchy--Schwarz estimate of it returns
exactly zero of the rung on every rung of the tested ladder. As an application we
compose the certified rungs with a margin-free bordering criterion into a chain that
transports semidefiniteness along a nested family of windows, and we report a verified
instance of that chain in which every load-bearing step is a completed Cholesky
factorisation above an explicitly declared floating-point floor: $430$ such
certificates over $52$ windows and $88$ rungs. The chain is verified window by window
and is not claimed uniformly in the window index. The example class that motivates the
construction has a symbol built from a logarithmic profile minus an explicit cosine
comb at the frequencies $\log n$, $n$ a prime power; the arithmetic enters only as the
source of the comb positions and only through two classical gap inequalities.
\end{abstract}

\section{Introduction}
\label{sec:intro}

\subsection{The problem, and why the direction matters}

Let $A\in\R^{h\times h}$ be symmetric positive definite, let $b\in\R^{h}$, and set
\begin{equation}
\label{eq:eps-def}
  \varepsilon \;=\; 1-b^{\T}A^{-1}b,
  \qquad
  Q \;=\; A-bb^{\T}.
\end{equation}
The scalar $\varepsilon$ is the value of the quadratic functional
$J(v)=1-2b^{\T}v+v^{\T}Av$ at its minimiser $u=A^{-1}b$; when $A$ is a Galerkin
matrix on a finite-dimensional space $V$ and $b$ the load vector of a fixed
right-hand side, $\varepsilon$ is exactly the Galerkin error of the best
approximation in $V$, measured in the energy of the ambient form. When $A$ is a
Toeplitz matrix and $b$ the first coordinate vector, $\varepsilon$ is the classical
one-step prediction error of Szeg\H{o}, Levinson and Durbin
\cite{Szego1939,GrenanderSzego1958,Levinson1947,Durbin1960}, whence the name we use
throughout. Finally, and this is the elementary fact the whole paper turns on, by
Albert's criterion applied to the two Schur readings of the bordered matrix
$\bigl[\begin{smallmatrix}A&b\\b^{\T}&1\end{smallmatrix}\bigr]$ one has, for $A\PD$,
\begin{equation}
\label{eq:albert-eps}
  \varepsilon\ge 0
  \quad\Longleftrightarrow\quad
  \begin{bmatrix}A & b\\ b^{\T} & 1\end{bmatrix}\PSD
  \quad\Longleftrightarrow\quad
  Q\PSD .
\end{equation}
So a \emph{lower} bound on $\varepsilon$ is a statement about the semidefiniteness of
a rank-one-perturbed form, and vice versa.

Upper bounds for $\varepsilon$ are the classical direction. C\'ea's lemma, duality,
and the Bramble--Hilbert machinery all estimate an approximation error from above
\cite{Cea1964}. Lower bounds are the hard direction, and the standard device for
obtaining them --- the \emph{saturation assumption}, which asserts that refining the
space reduces the error by a factor bounded away from one --- is known to be delicate:
Bank and Smith made it the basis of hierarchical a posteriori estimators
\cite{BankSmith1993}, Bornemann, Erdmann and Kornhuber isolated the constants it
consumes \cite{BornemannErdmannKornhuber1996}, and D\"orfler and Nochetto showed that
it holds under a smallness condition on the data oscillation and can fail otherwise
\cite{DorflerNochetto2002}. In the multilevel literature the same difficulty appears
as the strengthened Cauchy--Schwarz constant $\gamma^{2}<1$ of hierarchical
splittings \cite{Yserentant1986,BPX1990,AxelssonVassilevski1989,Brandt1977}: it is a
\emph{worst-case} quantity, an upper bound on a coupling, and any bound that has to
supply it inherits its worst case.

There is a second, more specific obstruction in the Toeplitz setting, and it is what
makes the example class of Section~\ref{sec:example} genuinely hard rather than
merely technical. If the symbol $\sigma$ of the form is nonnegative, the Toeplitz
matrix $\Tm{M}(\sigma)$ is positive semidefinite, and conversely a pointwise minorant
$\ell\le\sigma$ transfers to an operator minorant $\Tm{M}(\ell)\preceq\Tm{M}(\sigma)$;
this is Parseval together with the Fej\'er--Riesz representation of nonnegative
trigonometric polynomials \cite{Fejer1916,Riesz1916}, and it is the standard route
from a symbol to a spectral bound \cite{KMS1953,GrenanderSzego1958}. The route is
one-sided in a way that matters here. For the forms we care about the smallest
eigenvalue is not produced by a small value of the symbol at some frequency; it is
produced by \emph{cancellation inside the quadratic form}, along a direction that is
smooth and lives in the coarse part of the space. No pointwise minorant of the symbol
can see such a direction, because a pointwise minorant is by construction blind to
which vector realises the minimum. Section~\ref{sec:example} records the measured
form of this statement for the example class; the theorem of
Section~\ref{sec:main} is designed so that it never needs a symbol minorant at all.

\subsection{Contribution}

The result of the paper is (C1). The rest is what (C1) is good for, and what it costs.

\begin{enumerate}[label=(C\arabic*),leftmargin=3.1em]
\item \textbf{The result: a direction reversal that makes a certified rung possible}
  (Theorem~\ref{thm:rung}). On a nested ladder, the increment
  $\delta_{l}=\varepsilon_{l}-\varepsilon_{l+1}$ of the saturation telescope is the
  squared \emph{inverse}-energy norm of the level residual, hence a maximum over test
  vectors. A single test vector therefore supplies a valid lower bound, and the
  denominator of that bound wants the relevant Schur complement from \emph{above}.
  Schur complement monotonicity \cite{Haynsworth1968} supplies the upper bound from
  the mere semidefiniteness of the coarse block, Parseval supplies an upper bound by
  a certified majorant of the symbol, and L\"owner's antitony of the matrix inverse
  \cite{Loewner1934} closes the chain. No coarse inverse, no $\lmin$ of a coarse
  block, and no strengthened Cauchy--Schwarz constant enters. The optimal test vector
  for the majorant is available in closed form, and the entire loss relative to
  $\delta_{l}$ is the loss in the majorant. The nesting the theorem needs is not
  assumed: Lemma~\ref{lem:nesting} proves it for every lag family generated by a
  fixed kernel, and Lemma~\ref{lem:classconsistent} exhibits the kernel for the
  example class.
\item \textbf{The scope of the result, stated as such.} The numerator of the certified
  rung is evaluated at the coarse Galerkin solution carried along the ladder
  (Convention~\ref{conv:carried}), which is exactly the information an a posteriori
  estimator has and no more: it involves neither the fine solution nor
  $\varepsilon$ itself. Remark~\ref{rem:scope} spells this out, and
  Proposition~\ref{prop:negative} measures the distance to the stronger, a priori
  object: the solution-free variant of the same rung fails through cancellation, which
  we report as a negative result on a specific ladder because it is what tells a reader
  that the convention was not an oversight.
\item \textbf{An application: margin-free composition into a chain}
  (Corollary~\ref{cor:chain}), together with a verified instance of it. Two exact
  identities (Lemmas~\ref{lem:zeroblock} and \ref{lem:handover}) show that on a nested
  family of windows at one common resolution the trailing block of the next window's
  form \emph{is} the current window's form. Albert's criterion \cite{Albert1969} with
  Douglas' range condition \cite{Douglas1966} then transports semidefiniteness from
  one window to the next using only the \emph{sign} of that block --- never the size
  of $\varepsilon$ --- and the residual condition is a Schur complement of the small
  dimension of the window increment. The hypotheses of that corollary are verified
  window by window on the instance of Section~\ref{sec:numerics}; a version uniform in
  the window index is not claimed, and item~\ref{L:uniform} of
  Section~\ref{sec:limits} says precisely what is missing for one.
\end{enumerate}

The instance is produced in Section~\ref{sec:numerics} by a single script, runtime
about ten seconds: $430$ completed Cholesky certificates and the identity residuals,
each stated relative to a declared floating-point floor in the sense of Wilkinson and
Higham \cite{Wilkinson1963,Higham2002}. We are careful throughout to distinguish what
is proved (the identities and Theorem~\ref{thm:rung}) from what is certified on one
finite instance (the hypotheses of Corollary~\ref{cor:chain}) from what is merely
measured or fitted.

\subsection{Notation}

$\Tm{M}(c)=\bigl(c_{|r-s|}\bigr)_{r,s=0}^{M-1}$ is the symmetric Toeplitz matrix of a
real sequence $c=(c_{m})_{m\ge0}$. For real integrable even $g$ on $[-\pi,\pi]$ with
Fourier coefficients $\widehat g_{m}$ we write $\Tm{M}(g):=\Tm{M}(\widehat g)$;
no confusion arises because a symbol enters only through its first $M$ coefficients.
$X\PSD$ and $X\PD$ denote positive semidefiniteness and definiteness, $X\preceq Y$
means $Y-X\PSD$, $X^{+}$ is the Moore--Penrose pseudoinverse, and
$\|v\|_{X}^{2}=v^{\T}Xv$ for $X\PSD$. Finally
$(t)_{+}=\max(t,0)$ and $\Delta_{D}(s)=(1-|s|/D)_{+}$ is the unit-height hat of
half-width $D$.

\section{Setting}
\label{sec:setting}

\subsection{The odd window form and its defect}

Fix a cell width $D>0$ and an even $M\in2\N$, and put $\alpha=MD/2$ and $h=M/2$.
The grid is the partition of $(-\alpha,\alpha)$ into $M$ cells of width $D$ with
midpoints
\begin{equation}
\label{eq:cells}
  \bar x_{r} \;=\; -\alpha+\bigl(r+\tfrac12\bigr)D, \qquad r=0,\dots,M-1 ,
\end{equation}
and $\chi_{r}$ denotes the $L^{2}$-normalised indicator of cell $r$, i.e.\
$\chi_{r}=D^{-1/2}\mathbf 1_{[\bar x_{r}-D/2,\;\bar x_{r}+D/2]}$. Let
$J\in\R^{M\times M}$ be the reflection $(Jv)_{r}=v_{M-1-r}$ and let
\begin{equation}
\label{eq:B}
  B\in\R^{M\times h},\qquad Be_{r}=\tfrac{1}{\sqrt2}\,(e_{r}-e_{M-1-r}),
  \qquad r=0,\dots,h-1,
\end{equation}
be the isometry onto the reflection-odd sector, so $B^{\T}B=I_{h}$ and
$JB=-B$.

\begin{definition}[window form and pole vector]
\label{def:window}
Given a real sequence $c=(c_{m})_{m=0}^{M-1}$, the \emph{odd window form} and the
\emph{pole vector} at $(D,M)$ are
\begin{align}
\label{eq:Adef}
  A &\;=\; B^{\T}\,\Tm{M}(c)\,B \;\in\;\R^{h\times h},
  &&A_{rs}=c_{|r-s|}-c_{M-1-r-s},\\
\label{eq:bdef}
  b &\;=\; \bigl(b_{r}\bigr)_{r=0}^{h-1},
  &&b_{r}=\frac{8}{\sqrt D}\,\sinh\!\Bigl(\frac D4\Bigr)\sinh\!\Bigl(\frac{\bar x_{r}}{2}\Bigr).
\end{align}
The associated \emph{defect} and \emph{bordered form} are $\varepsilon$ and $Q$ of
\eqref{eq:eps-def}.
\end{definition}

The particular constant in \eqref{eq:bdef} is not a normalisation choice; it is the
reason the whole construction is a Galerkin problem on nested spaces.

\begin{lemma}[$b$ is the coefficient vector of a fixed function]
\label{lem:bcoeff}
Let $f(x)=2\sinh(x/2)$. Then $b_{r}=\langle f,\chi_{r}\rangle_{L^{2}(\R)}$ for
$r=0,\dots,h-1$, for every $D>0$ and every $M$.
\end{lemma}

\begin{proof}
$\langle f,\chi_{r}\rangle
 =D^{-1/2}\int_{\bar x_{r}-D/2}^{\bar x_{r}+D/2}2\sinh(x/2)\,dx
 =D^{-1/2}\bigl[4\cosh(x/2)\bigr]_{\bar x_{r}-D/2}^{\bar x_{r}+D/2}$,
and $\cosh(a+d)-\cosh(a-d)=2\sinh a\sinh d$ with $a=\bar x_{r}/2$, $d=D/4$ gives
$8D^{-1/2}\sinh(D/4)\sinh(\bar x_{r}/2)$.
\end{proof}

\begin{remark}
\label{rem:galerkin}
Write $V_{D,M}\subset L^{2}(-\alpha,\alpha)$ for the space of odd piecewise-constant
functions on the grid \eqref{eq:cells}; $\{B\chi\}$ is an orthonormal basis of it, and
$\dim V_{D,M}=h$. If $a(\cdot,\cdot)$ denotes the bilinear form whose matrix in that
basis is $A$, then by Lemma~\ref{lem:bcoeff}
\begin{equation}
\label{eq:eps-galerkin}
  \varepsilon \;=\; \min_{v\in V_{D,M}}\Bigl(1-2\langle f,v\rangle+a(v,v)\Bigr),
\end{equation}
so $\varepsilon$ is the Galerkin defect of the fixed function $f$ in the space
$V_{D,M}$, and its minimiser is the Galerkin solution $u=A^{-1}b$. Nothing in
\eqref{eq:eps-galerkin} refers to $c$ except through $a$.
\end{remark}

\subsection{The ladder and exact nesting}
\label{sec:ladder}

We refine at fixed $\alpha$. Let $L\ge1$ and
\begin{equation}
\label{eq:levels}
  D_{l}=2^{-l}D_{0},\qquad M_{l}=2^{l}M_{0},\qquad h_{l}=M_{l}/2,
  \qquad l=0,\dots,L,
\end{equation}
so $\alpha=M_{l}D_{l}/2$ is the same at every level and
$V_{D_{0},M_{0}}\subset V_{D_{1},M_{1}}\subset\dots\subset V_{D_{L},M_{L}}$. Write
$A^{(l)},b^{(l)},\varepsilon_{l},u^{(l)}$ for the objects of
Definition~\ref{def:window} at level $l$. Let $P_{l}:\R^{h_{l}}\to\R^{h_{l+1}}$ be
the matrix of the inclusion $V_{D_{l},M_{l}}\hookrightarrow V_{D_{l+1},M_{l+1}}$ in
the orthonormal bases, i.e.\ the $L^{2}$-normalised piecewise-constant prolongation:
each coarse cell splits into two fine cells and
\begin{equation}
\label{eq:P}
  (P_{l}x)_{2i}=(P_{l}x)_{2i+1}=\tfrac{1}{\sqrt2}\,x_{i}.
\end{equation}
We suppress the level index on $P$ when it is clear. Let
$Z_{l}:\R^{h_{l}}\to\R^{h_{l+1}}$ be the complementary Haar injection,
\begin{equation}
\label{eq:Z}
  (Z_{l}x)_{2i}=\tfrac{1}{\sqrt2}\,x_{i},\qquad
  (Z_{l}x)_{2i+1}=-\tfrac{1}{\sqrt2}\,x_{i},
\end{equation}
so that
\begin{equation}
\label{eq:isometries}
  P^{\T}P=Z^{\T}Z=I_{h_{l}},\qquad P^{\T}Z=0,\qquad
  \ran P\oplus\ran Z=\R^{h_{l+1}} .
\end{equation}
The pair $(P,Z)$ is a two-level orthogonal splitting of the fine space into its
piecewise-constant (coarse) and oscillatory (high-pass) parts.

The property we need of the ladder is that it is \emph{one} bilinear form on nested
spaces:
\begin{equation}
\label{eq:nesting}
  P^{\T}A^{(l+1)}P=A^{(l)},
  \qquad
  P^{\T}b^{(l+1)}=b^{(l)}
  \qquad (l<L).
\end{equation}
Both halves of \eqref{eq:nesting} are identities, and we prove them. The second is
unconditional.

\begin{lemma}[nesting of the pole vector]
\label{lem:bnesting}
$P^{\T}b^{(l+1)}=b^{(l)}$ holds for every $D_{l}$ and every $M_{l}$.
\end{lemma}

\begin{proof}
Immediate from Lemma~\ref{lem:bcoeff}, since $P$ is the matrix of an inclusion of
spaces in orthonormal bases and $f$ is fixed. Directly: with $d=D_{l}/2$ the two fine
midpoints inside the coarse cell centred at $\bar x$ are $\bar x\mp d/2$, so
\[
  (P^{\T}b^{(l+1)})_{i}
  =\frac{1}{\sqrt2}\cdot\frac{8}{\sqrt d}\sinh\frac d4
   \Bigl[\sinh\frac{\bar x-d/2}{2}+\sinh\frac{\bar x+d/2}{2}\Bigr]
  =\frac{1}{\sqrt2}\cdot\frac{8}{\sqrt d}\sinh\frac d4\cdot
   2\cosh\frac d4\,\sinh\frac{\bar x}{2},
\]
which equals $8\,(2d)^{-1/2}\sinh(d/2)\sinh(\bar x/2)=b^{(l)}_{i}$ because
$D_{l}=2d$.
\end{proof}

The first half holds for every lag family generated by a fixed kernel, in the
following sense.

\begin{definition}[consistent lag family]
\label{def:consistent}
Let $K$ be an even distribution on $\R$ for which $K*\Delta_{D}$ is a well-defined
continuous function for every $D>0$. The family $\{c^{(l)}\}_{l=0}^{L}$ is
\emph{generated by} $K$ if
\begin{equation}
\label{eq:consistent}
  c^{(l)}_{m}\;=\;\bigl(K*\Delta_{D_{l}}\bigr)(mD_{l}),
  \qquad 0\le m\le M_{l}-1,\quad 0\le l\le L,
\end{equation}
with the \emph{same} $K$ at every level. We then call the family \emph{consistent}.
\end{definition}

Condition \eqref{eq:consistent} is not an extra hypothesis dressed up as a
definition; it is what the phrase ``the Toeplitz matrix of a kernel'' means. Indeed if
$T_{rs}=\iint K(x-y)\chi_{r}(x)\chi_{s}(y)\,dx\,dy$ is the Galerkin matrix of
convolution by $K$ in the cell basis, then substituting $t=y-x$ and using that the
convolution of two indicators of width $D$ shifted by $mD$ is $D\,\Delta_{D}(\cdot-mD)$
gives $T_{rs}=(K*\Delta_{D})\bigl((r-s)D\bigr)$, i.e.\ exactly \eqref{eq:consistent}.
The even extension $c^{(l)}_{-m}:=c^{(l)}_{m}$ is automatic, because $K$ even implies
$K*\Delta_{D}$ even.

\begin{lemma}[hat refinement]
\label{lem:hat}
For every $d>0$ and every $u\in\R$,
\begin{equation}
\label{eq:hatref}
  \Delta_{d}(u)+\tfrac12\Delta_{d}(u-d)+\tfrac12\Delta_{d}(u+d)
  \;=\;\Delta_{2d}(u).
\end{equation}
\end{lemma}

\begin{proof}
Both sides are continuous, piecewise linear, supported in $[-2d,2d]$, and their
breakpoints are contained in $\{-2d,-d,0,d,2d\}$: the left-hand side because each of
its three terms is piecewise linear with breakpoints in that set, the right-hand side
because $\Delta_{2d}$ has breakpoints $\{-2d,0,2d\}$. Two continuous piecewise-linear
functions with breakpoints in a common finite set agree everywhere as soon as they
agree at every point of that set. At $u=0$ the left-hand side is
$1+\tfrac12\cdot0+\tfrac12\cdot0=1=\Delta_{2d}(0)$; at $u=\pm d$ it is
$0+\tfrac12\cdot1+0=\tfrac12=\Delta_{2d}(\pm d)$; at $u=\pm2d$ both sides vanish.
\end{proof}

\begin{lemma}[exact nesting]
\label{lem:nesting}
Let $\{c^{(l)}\}$ be consistent in the sense of Definition~\ref{def:consistent}. Then
\eqref{eq:nesting} holds for every $l<L$.
\end{lemma}

\begin{proof}
\emph{Step 1: reduction to the full Toeplitz level.} Let
$P_{\mathrm{full}}:\R^{M_{l}}\to\R^{M_{l+1}}$ be the prolongation \eqref{eq:P} on the
unreduced grid. Then $B_{M_{l+1}}P=P_{\mathrm{full}}B_{M_{l}}$. Indeed, using
$M_{l+1}=2M_{l}$,
\begin{align*}
  P_{\mathrm{full}}B_{M_{l}}e_{r}
  &=\tfrac{1}{\sqrt2}P_{\mathrm{full}}\bigl(e_{r}-e_{M_{l}-1-r}\bigr)\\
  &=\tfrac12\Bigl[\bigl(e_{2r}-e_{M_{l+1}-1-2r}\bigr)
      +\bigl(e_{2r+1}-e_{M_{l+1}-1-(2r+1)}\bigr)\Bigr]
  =B_{M_{l+1}}\tfrac{e_{2r}+e_{2r+1}}{\sqrt2},
\end{align*}
because $2(M_{l}-1-r)=M_{l+1}-1-(2r+1)$ and $2(M_{l}-1-r)+1=M_{l+1}-1-2r$; and the
last expression is $B_{M_{l+1}}Pe_{r}$ by \eqref{eq:P}. Consequently
\begin{equation}
\label{eq:reduction}
  P^{\T}A^{(l+1)}P
  =\bigl(B_{M_{l+1}}P\bigr)^{\T}\Tm{M_{l+1}}(c^{(l+1)})\bigl(B_{M_{l+1}}P\bigr)
  =B_{M_{l}}^{\T}\Bigl(P_{\mathrm{full}}^{\T}\Tm{M_{l+1}}(c^{(l+1)})
       P_{\mathrm{full}}\Bigr)B_{M_{l}},
\end{equation}
so it suffices to prove
$P_{\mathrm{full}}^{\T}\Tm{M_{l+1}}(c^{(l+1)})P_{\mathrm{full}}=\Tm{M_{l}}(c^{(l)})$.

\emph{Step 2: the coarsening is a three-term average.} Writing $c'=c^{(l+1)}$ and
using that the only nonzero entries of $P_{\mathrm{full}}$ in column $i$ are the two
rows $2i,2i+1$, each equal to $2^{-1/2}$,
\[
  \bigl(P_{\mathrm{full}}^{\T}\Tm{M_{l+1}}(c')P_{\mathrm{full}}\bigr)_{ik}
  =\tfrac12\!\!\sum_{a\in\{2i,2i+1\}}\;\sum_{b\in\{2k,2k+1\}}\!\! c'_{|a-b|}
  =c'_{2|m|}+\tfrac12\bigl(c'_{|2m-1|}+c'_{|2m+1|}\bigr),
  \qquad m=i-k,
\]
since the four terms are $c'_{|2m|}$, $c'_{|2m-1|}$, $c'_{|2m+1|}$ and $c'_{|2m|}$.
The largest index occurring is $2(M_{l}-1)+1=M_{l+1}-1$, so every entry used is
available. The matrix on the left is therefore Toeplitz, and
$P_{\mathrm{full}}^{\T}\Tm{M_{l+1}}(c')P_{\mathrm{full}}=\Tm{M_{l}}(\tilde c\,)$ with
\begin{equation}
\label{eq:threeterm}
  \tilde c_{n}=c'_{2n}+\tfrac12\bigl(c'_{2n-1}+c'_{2n+1}\bigr),
  \qquad 0\le n\le M_{l}-1,
\end{equation}
under the even extension $c'_{-1}:=c'_{1}$.

\emph{Step 3: consistency turns \eqref{eq:threeterm} into $\tilde c=c^{(l)}$.} Put
$d=D_{l+1}=D_{l}/2$. By \eqref{eq:consistent} and $(K*g)(x+a)=(K*g(\cdot+a))(x)$,
\[
  \tilde c_{n}
  =\bigl(K*\Delta_{d}\bigr)(2nd)
   +\tfrac12\bigl(K*\Delta_{d}\bigr)\bigl(2nd-d\bigr)
   +\tfrac12\bigl(K*\Delta_{d}\bigr)\bigl(2nd+d\bigr)
  =\bigl(K*G_{d}\bigr)(2nd),
\]
where $G_{d}(u)=\Delta_{d}(u)+\tfrac12\Delta_{d}(u+d)+\tfrac12\Delta_{d}(u-d)$. By
Lemma~\ref{lem:hat}, $G_{d}=\Delta_{2d}=\Delta_{D_{l}}$, whence
$\tilde c_{n}=(K*\Delta_{D_{l}})(nD_{l})=c^{(l)}_{n}$. Together with
\eqref{eq:reduction} this is the first half of \eqref{eq:nesting}; the second half is
Lemma~\ref{lem:bnesting}.
\end{proof}

The example class of Section~\ref{sec:example} is consistent, with an explicit $K$
(Lemma~\ref{lem:classconsistent}), so \eqref{eq:nesting} is an identity there and not
a hypothesis. What Section~\ref{sec:numerics} reports for it is a floating-point
statement about two different orders of summation for the same exact quantity: the
matrix $P^{\T}A^{(l+1)}P$, assembled at the fine level and then coarsened, agrees with
$A^{(l)}$, assembled directly at the coarse level, to relative $3.0\cdot10^{-14}$.

\subsection{The two-level split of the fine form}
\label{sec:split}

Fix a level $l$ and abbreviate $A_{f}=A^{(l+1)}$, $V=P$, $W=Z$. Setting
\begin{equation}
\label{eq:blocks}
  A_{c}=V^{\T}A_{f}V,\qquad
  A_{z}=W^{\T}A_{f}W,\qquad
  B_{x}=V^{\T}A_{f}W,
\end{equation}
the orthogonal splitting \eqref{eq:isometries} puts the fine form in block shape
$\bigl[\begin{smallmatrix}A_{c}&B_{x}\\ B_{x}^{\T}&A_{z}\end{smallmatrix}\bigr]$, and
the nesting \eqref{eq:nesting} identifies $A_{c}=A^{(l)}$: \emph{the coarse block of a
rung is the window form of the level below it}. The Schur complement onto the
oscillatory space is
\begin{equation}
\label{eq:Sdef}
  S \;=\; A_{z}-B_{x}^{\T}A_{c}^{-1}B_{x} .
\end{equation}
If $A_{f}\PD$ then $A_{c}=V^{\T}A_{f}V\PD$ for the isometry $V$, so \eqref{eq:Sdef}
is well defined, and $S\PD$.

\subsection{The example class}
\label{sec:example}

The construction above is symbol-agnostic; Theorem~\ref{thm:rung} and
Corollary~\ref{cor:chain} use no property of $c$ beyond what is stated in their
hypotheses. The class we verify on, and which motivates the whole question, is the
following.

\begin{definition}[the comb class]
\label{def:class}
Let $a:\R\to\R$ be the even kernel
\begin{equation}
\label{eq:arch}
  a(s)=-(\gamma+\log\pi)\,\Delta_{D}(s)
      -\Delta_{D}(s)\log\!\bigl(1-e^{-2W}\bigr)
      +2\int_{0}^{W}
        \frac{\Delta_{D}(s)e^{-2w}-S_{D}(s,w)\,e^{-w/2}}{1-e^{-2w}}\,dw,
\end{equation}
where $W=|s|+D$, $\;S_{D}(s,w)=\tfrac12\bigl[\Delta_{D}(s-w)+\Delta_{D}(s+w)\bigr]$,
and $\gamma$ is Euler's constant. Let $\mathcal U=\{(u_{j},\mu_{j})\}_{j}$ be a
finite set of \emph{comb nodes} $u_{j}>0$ with weights $\mu_{j}>0$, and put
\begin{equation}
\label{eq:lags}
  c_{m} \;=\; a(mD)\;-\;\sum_{j:\,u_{j}\le 2\alpha}\mu_{j}\,S_{D}(mD,u_{j}),
  \qquad m=0,\dots,M-1 .
\end{equation}
The \emph{prime-power instance} is $u_{j}=\log n_{j}$ over the prime powers
$n_{j}=p^{\nu}$ in increasing order, with $\mu_{j}=2\Lambda(n_{j})/\sqrt{n_{j}}$ and
$\Lambda$ the von Mangoldt weight $\Lambda(p^{\nu})=\log p$.
\end{definition}

The form \eqref{eq:arch} is written the way it is computed: the upper limit $W=|s|+D$
truncates the integral at the last point where the second numerator can be nonzero,
and the $\log(1-e^{-2W})$ term is the exact remainder of the tail. Undoing that
bookkeeping identifies the fixed kernel behind the class.

\begin{lemma}[the comb class is consistent]
\label{lem:classconsistent}
Fix $\alpha$ and a node set $\mathcal U$, and let $c^{(l)}$ be given by
\eqref{eq:lags} at $(D_{l},M_{l})$ along the ladder \eqref{eq:levels}. Then
$\{c^{(l)}\}$ is consistent in the sense of Definition~\ref{def:consistent}, generated
by the even distribution
\begin{equation}
\label{eq:K}
  K=\underbrace{-(\gamma+\log\pi)\,\delta_{0}
    +2\int_{0}^{\infty}\frac{e^{-2w}\delta_{0}
       -e^{-w/2}\,\tfrac12(\delta_{w}+\delta_{-w})}{1-e^{-2w}}\,dw}_{K_{\mathrm{arch}}}
  \;-\;\sum_{j:\,u_{j}\le2\alpha}\mu_{j}\,
       \tfrac12\bigl(\delta_{u_{j}}+\delta_{-u_{j}}\bigr),
\end{equation}
which does not depend on $l$. Consequently Lemma~\ref{lem:nesting} applies and
\eqref{eq:nesting} holds identically on the ladder.
\end{lemma}

\begin{proof}
Both parts are pairings of a fixed even distribution against the hat, since
$\bigl(\delta_{v}*\Delta_{D}\bigr)(s)=\Delta_{D}(s-v)$ and hence
$\bigl(\tfrac12(\delta_{v}+\delta_{-v})*\Delta_{D}\bigr)(s)=S_{D}(s,v)$.

\emph{The comb part.} The node set in \eqref{eq:lags} is $\{u_{j}\le2\alpha\}$, and
$\alpha=M_{l}D_{l}/2$ is the \emph{same} at every level of the ladder
\eqref{eq:levels}; so the sum in \eqref{eq:lags} runs over one and the same finite set
for all $l$, and it equals $\bigl(\sum_{j}\mu_{j}\tfrac12(\delta_{u_{j}}
+\delta_{-u_{j}})*\Delta_{D_{l}}\bigr)(mD_{l})$ by the displayed identity.

\emph{The archimedean part.} Fix $s$ and let $w>W=|s|+D$. Then
$\Delta_{D}(s-w)=0$, because $w-s\ge w-|s|>D$, and $\Delta_{D}(s+w)=0$, because
$s+w\ge w-|s|>D$; hence $S_{D}(s,w)=0$ and the integrand of \eqref{eq:arch} reduces to
$\Delta_{D}(s)e^{-2w}/(1-e^{-2w})$ there. Since
$\frac{d}{dw}\log(1-e^{-2w})=2e^{-2w}/(1-e^{-2w})$,
\[
  2\int_{W}^{\infty}\frac{\Delta_{D}(s)e^{-2w}}{1-e^{-2w}}\,dw
  =\Delta_{D}(s)\Bigl[\log\bigl(1-e^{-2w}\bigr)\Bigr]_{W}^{\infty}
  =-\Delta_{D}(s)\log\bigl(1-e^{-2W}\bigr),
\]
which is exactly the second term of \eqref{eq:arch}. Therefore
\begin{equation}
\label{eq:archinfty}
  a(s)=-(\gamma+\log\pi)\Delta_{D}(s)
   +2\int_{0}^{\infty}
     \frac{\Delta_{D}(s)e^{-2w}-S_{D}(s,w)\,e^{-w/2}}{1-e^{-2w}}\,dw ,
\end{equation}
an integral that converges absolutely at $w=0$ because the numerator vanishes there
($S_{D}(s,0)=\Delta_{D}(s)$) and is $O(w)$, while the denominator is $\sim2w$; and at
$w=\infty$ because of $e^{-w/2}$. Reading \eqref{eq:archinfty} through the displayed
identity gives $a=K_{\mathrm{arch}}*\Delta_{D}$ with $K_{\mathrm{arch}}$ as in
\eqref{eq:K}, independent of $D$. Adding the two parts and sampling at $s=mD_{l}$
gives \eqref{eq:consistent}.
\end{proof}

\begin{remark}
\label{rem:noassumption}
Lemma~\ref{lem:classconsistent} is the reason this paper carries no nesting
assumption. It also isolates what the ladder \eqref{eq:levels} is for: refining at
\emph{fixed} $\alpha$ keeps the admitted node set constant, and a family whose node
set changed with the level would not be consistent.
\end{remark}

Three comments on the class, all of them limitations of scope rather than claims.

\emph{(i) The label ``log-singular''.} The integrand of \eqref{eq:arch} carries the
factor $(1-e^{-2w})^{-1}=\sum_{n\ge0}e^{-2nw}$, whose expansion is the digamma
series, and the constant is $-\gamma-\log\pi$. The kernel is therefore a
logarithmic (digamma) profile smoothed against the hat $\Delta_{D}$, and the label
in the title refers to that profile. The integrand vanishes at $w=0$ because
$S_{D}(s,0)=\Delta_{D}(s)$, so $a$ is finite for every $s$; the symbol
$\sigma^{(M)}(\theta)=\sum_{|m|\le M-1}c_{|m|}e^{im\theta}$ is a trigonometric
polynomial for each fixed $M$, and no asymptotic property of the family is used
anywhere below. The classical regularity theory for symbols of this type is Widom's
and Fisher--Hartwig's \cite{Widom1974,FisherHartwig1969,BoettcherSilbermann2006}; we
name it as the address of the questions we do \emph{not} answer, not as an ingredient.

\emph{(ii) What the comb does.} By \eqref{eq:lags} the symbol of $c$ is a logarithmic
profile minus an explicit cosine comb: each node contributes a term proportional to
$\cos(u_{j}\theta/D)$ modulated by the hat's Fourier transform. This is the structural
reason a pointwise symbol minorant is useless here. The comb dips are numerous and
narrow --- narrow enough that the large-sieve budget one would naturally spend on them
\cite{MontgomeryVaughan1974} is exhausted before the window is covered --- the
infimum of the symbol over a window is driven to zero by them, and the
near-null direction of the resulting form is nevertheless \emph{smooth}: it lives in
the coarse part $\ran V$ and its Rayleigh quotient is small because of cancellation
inside the form, not because the symbol is small along it. The measured statement is
in Section~\ref{sec:numerics}.

\emph{(iii) What the arithmetic is used for.} Only the positions $u_{j}=\log n_{j}$,
and only through two classical inequalities on the consecutive log-gaps
$g_{k}=u_{k+1}-u_{k}$, namely $g_{k}\le\log 2$ (Bertrand's postulate in Chebyshev's
form \cite{Chebyshev1852}) and $g_{k}\ge\log(1+1/n_{k})$. Both are used solely to
check that a resolution is admissible in the sense of
Definition~\ref{def:admissible}; no unproved property of the sequence
$(u_{j})$ enters the construction, and no arithmetic statement is claimed.

\section{The saturation telescope}
\label{sec:telescope}

The identities of this section are classical; we record them in the form we use.

\begin{lemma}[saturation identity and telescope]
\label{lem:telescope}
Assume $A^{(l)}\PD$ for all $l$ and the nesting \eqref{eq:nesting}. Put
\begin{equation}
\label{eq:delta}
  \delta_{l}\;=\;\bigl\|u^{(l+1)}-Pu^{(l)}\bigr\|_{A^{(l+1)}}^{2}\;\ge\;0 .
\end{equation}
Then for every $l<L$
\begin{equation}
\label{eq:sat}
  \varepsilon_{l}\;=\;\varepsilon_{l+1}+\delta_{l},
  \qquad\text{hence}\qquad
  \varepsilon_{0}\;=\;\varepsilon_{L}+\sum_{l=0}^{L-1}\delta_{l}
  \quad\text{(exactly)} .
\end{equation}
Moreover, with the level residual $r_{l}=b^{(l+1)}-A^{(l+1)}Pu^{(l)}$,
\begin{equation}
\label{eq:dual}
  V^{\T}r_{l}=0,
  \qquad
  \delta_{l}=\bigl\|r_{l}\bigr\|^{2}_{(A^{(l+1)})^{-1}} .
\end{equation}
\end{lemma}

\begin{proof}
Let $J_{l+1}(v)=1-2(b^{(l+1)})^{\T}v+v^{\T}A^{(l+1)}v$ on $\R^{h_{l+1}}$. Completing
the square at its minimiser $u^{(l+1)}$ gives
$J_{l+1}(v)=\varepsilon_{l+1}+\|v-u^{(l+1)}\|^{2}_{A^{(l+1)}}$ for all $v$. Take
$v=Pu^{(l)}$. By \eqref{eq:nesting},
$J_{l+1}(Pu^{(l)})=1-2(b^{(l)})^{\T}u^{(l)}+(u^{(l)})^{\T}A^{(l)}u^{(l)}
=\varepsilon_{l}$, which is \eqref{eq:sat}; the telescope follows by summation. For
\eqref{eq:dual}, \eqref{eq:nesting} gives
$V^{\T}r_{l}=b^{(l)}-A^{(l)}u^{(l)}=0$, which is Galerkin orthogonality, and
$A^{(l+1)}(u^{(l+1)}-Pu^{(l)})=r_{l}$ turns \eqref{eq:delta} into \eqref{eq:dual}.
\end{proof}

Two consequences of \eqref{eq:dual} deserve to be separated, because the whole
difficulty of the subject sits in the difference between them.

\begin{lemma}[the two dual descriptions of a rung]
\label{lem:twoforms}
In the notation of Section~\ref{sec:split}, put
\begin{equation}
\label{eq:shat}
  s_{l}=W^{\T}b^{(l+1)},\qquad
  \shat_{l}=s_{l}-B_{x}^{\T}u^{(l)}\;=\;W^{\T}r_{l} .
\end{equation}
Then
\begin{equation}
\label{eq:primal}
  \delta_{l}
  \;=\;\min_{x\in\R^{h_{l}}}\;
    \bigl(Vx+W y^{\star}\bigr)^{\T}A^{(l+1)}\bigl(Vx+Wy^{\star}\bigr)
  \quad\text{for a suitable }y^{\star},
\end{equation}
a \emph{minimum}, and
\begin{equation}
\label{eq:maximum}
  \delta_{l}
  \;=\;\shat_{l}^{\T}S^{-1}\shat_{l}
  \;=\;\max_{v\neq0}\frac{(\shat_{l}^{\T}v)^{2}}{v^{\T}Sv},
\end{equation}
a \emph{maximum}.
\end{lemma}

\begin{proof}
$\delta_{l}=\|u^{(l+1)}-Pu^{(l)}\|^{2}_{A^{(l+1)}}$ is the squared energy distance
from $u^{(l+1)}$ to the coarse space, which is \eqref{eq:primal}. For
\eqref{eq:maximum}, \eqref{eq:dual} gives $r_{l}\in\ran W$, so $r_{l}=W\shat_{l}$ by
\eqref{eq:isometries}, and
$\delta_{l}=r_{l}^{\T}(A^{(l+1)})^{-1}r_{l}
=\shat_{l}^{\T}\bigl(W^{\T}(A^{(l+1)})^{-1}W\bigr)\shat_{l}$.
The block inverse formula for the splitting \eqref{eq:blocks} gives
$W^{\T}(A^{(l+1)})^{-1}W=S^{-1}$. The variational identity
$\shat^{\T}S^{-1}\shat=\max_{v}(\shat^{\T}v)^{2}/(v^{\T}Sv)$ for $S\PD$ is
Cauchy--Schwarz in the $S$-inner product, with equality at $v=S^{-1}\shat$.
\end{proof}

\begin{remark}[the direction, stated pedantically]
\label{rem:direction}
From \eqref{eq:primal}, every test vector yields an \emph{upper} bound on
$\delta_{l}$; a minimum principle cannot supply a lower bound, no matter how good the
test vector. From \eqref{eq:maximum}, every test vector yields a \emph{lower} bound
on $\delta_{l}$, provided the denominator $v^{\T}Sv$ is replaced by something
\emph{larger}. This is the direction reversal of the title of
Section~\ref{sec:main}: what a lower bound on $\delta_{l}$ needs is an
\emph{upper} bound on $S$, and upper bounds on Schur complements are cheap.
\end{remark}

\section{The certified rung}
\label{sec:main}

We need two classical upper-bound mechanisms.

\begin{lemma}[Schur complement monotonicity]
\label{lem:haynsworth}
With the notation of \eqref{eq:blocks} and $A_{c}\PD$, one has $S\preceq A_{z}$, and
the only property of $A_{c}$ used is $A_{c}^{-1}\PSD$, i.e.\ the \emph{sign} of
$A_{c}$.
\end{lemma}

\begin{proof}
$A_{z}-S=B_{x}^{\T}A_{c}^{-1}B_{x}\PSD$ whenever $A_{c}^{-1}\PSD$. See
\cite{Haynsworth1968}.
\end{proof}

\begin{lemma}[Parseval majorant]
\label{lem:parseval}
Let $g_{1},g_{2}$ be real, even and integrable on $[-\pi,\pi]$ with
$g_{1}\le g_{2}$ almost everywhere. Then $\Tm{M}(g_{1})\preceq \Tm{M}(g_{2})$, and
consequently $R^{\T}\Tm{M}(g_{1})R\preceq R^{\T}\Tm{M}(g_{2})R$ for every real
matrix $R$ with $M$ rows.
\end{lemma}

\begin{proof}
For $v\in\R^{M}$ and $\mathcal V(\theta)=\sum_{r}v_{r}e^{ir\theta}$, Parseval gives
$v^{\T}\Tm{M}(g)v=\frac1{2\pi}\int_{-\pi}^{\pi}g(\theta)|\mathcal V(\theta)|^{2}\,d\theta$.
Apply this to $g=g_{2}-g_{1}\ge0$. (That every nonnegative trigonometric polynomial
is such a $|\mathcal V|^{2}$ is the Fej\'er--Riesz theorem
\cite{Fejer1916,Riesz1916}; only the displayed identity is needed here.)
\end{proof}

\begin{definition}[certified majorant]
\label{def:majorant}
Let $\sigma^{(M)}(\theta)=\sum_{|m|\le M-1}c_{|m|}e^{im\theta}$, so that
$\Tm{M}(c)=\Tm{M}(\sigma^{(M)})$. A real even
$\mathrm{up}\in L^{1}[-\pi,\pi]$ is a \emph{certified majorant} at level $l+1$ if
$\mathrm{up}\ge\sigma^{(M_{l+1})}$ pointwise. We then set
\begin{equation}
\label{eq:U}
  U_{l} \;=\; W^{\T}B^{\T}\,\Tm{M_{l+1}}(\mathrm{up})\,B\,W .
\end{equation}
\end{definition}

Note that $A_{z}=W^{\T}B^{\T}\Tm{M_{l+1}}(c)BW$ by \eqref{eq:Adef} and
\eqref{eq:blocks}, so Lemma~\ref{lem:parseval} applied with $R=BW$ gives
$A_{z}\preceq U_{l}$ directly, and the Hankel-type term in \eqref{eq:Adef} never has
to be estimated: it is the reflection half of the isometry $B$.

\begin{theorem}[the certified rung]
\label{thm:rung}
Fix a level $l$ and assume
\begin{enumerate}[label=\textup{(H\arabic*)},leftmargin=3.2em]
\item\label{H:pd} $A^{(l+1)}\PD$;
\item\label{H:nest} the nesting \eqref{eq:nesting} at level $l$ --- which by
      Lemma~\ref{lem:nesting} holds for every consistent lag family, in particular for
      the class of Definition~\ref{def:class} --- so that $A_{c}=A^{(l)}$;
\item\label{H:maj} $\mathrm{up}$ is a certified majorant at level $l+1$ and
      $U_{l}\PD$.
\end{enumerate}
Then, with $\shat_{l}$ of \eqref{eq:shat},
\begin{equation}
\label{eq:8R}
  \delta_{l} \;\ge\; \shat_{l}^{\T}U_{l}^{-1}\shat_{l} \;=:\; \widehat c_{l}
  \;\ge\;0 ,
\end{equation}
with equality in the second place only if $\shat_{l}=0$. The bound is attained by the
explicit test vector $v^{\star}=U_{l}^{-1}\shat_{l}$ in \eqref{eq:maximum} after
replacing $S$ by $U_{l}$. The proof uses no inverse, no smallest eigenvalue and no
condition number of $A_{c}$, and no strengthened Cauchy--Schwarz constant.
\end{theorem}

\begin{proof}
By \ref{H:pd} and Section~\ref{sec:split}, $A_{c}\PD$ and $S\PD$, so
Lemma~\ref{lem:twoforms} applies and $\delta_{l}=\shat_{l}^{\T}S^{-1}\shat_{l}$.
By Lemma~\ref{lem:haynsworth}, $S\preceq A_{z}$; here \ref{H:nest} is what makes
$A_{c}\PSD$ the conclusion of the level below rather than a new hypothesis. By
Definition~\ref{def:majorant} and Lemma~\ref{lem:parseval},
$A_{z}\preceq U_{l}$. Hence $0\prec S\preceq U_{l}$, and L\"owner's antitony of the
matrix inverse on the positive definite cone \cite{Loewner1934} gives
$U_{l}^{-1}\preceq S^{-1}$, i.e.
\[
  \delta_{l}=\shat_{l}^{\T}S^{-1}\shat_{l}
  \;\ge\;\shat_{l}^{\T}U_{l}^{-1}\shat_{l}=\widehat c_{l}.
\]
Positivity of $\widehat c_{l}$ for $\shat_{l}\neq0$ is $U_{l}\PD$. For the test
vector, substituting $v=U_{l}^{-1}\shat_{l}$ into
$(\shat_{l}^{\T}v)^{2}/(v^{\T}U_{l}v)$ returns $\shat_{l}^{\T}U_{l}^{-1}\shat_{l}$
exactly, so the closed-form vector is optimal for the majorant.
\end{proof}

\begin{remark}[what is certifiable, and how]
\label{rem:certify}
Each of \ref{H:pd}--\ref{H:maj} reduces to a completed Cholesky factorisation:
$A^{(l+1)}\PD$ and $U_{l}\PD$ directly, and the majorant property through
$\Tm{M_{l+1}}(\mathrm{up})-\Tm{M_{l+1}}(c)\PSD$, which is
Lemma~\ref{lem:parseval} in the certified direction. In particular
$A_{c}\PSD$ --- the hypothesis whose quantitative version blocks the classical
routes --- is supplied by the previous rung's own conclusion and never has to be
quantified. This is what makes a coarse-to-fine induction possible; see
Corollary~\ref{cor:chain}.
\end{remark}

\begin{remark}[the price of one test vector]
\label{rem:kantorovich}
The bound \eqref{eq:8R} is the exact value of the maximum \eqref{eq:maximum} with $S$
replaced by the majorant $U_{l}$, so the entire loss relative to $\delta_{l}$ is the
loss in $S\preceq U_{l}$ and none of it is a loss in the choice of test vector. The
classical estimate of what a single test vector costs against the optimal one
\cite{Kantorovich1948} is therefore quoted here only for the contrast: it does not
enter.
\end{remark}

\begin{convention}[carried-solution accounting]
\label{conv:carried}
The numerator $\shat_{l}=s_{l}-B_{x}^{\T}u^{(l)}$ of \eqref{eq:8R} is a
function of the coarse Galerkin solution $u^{(l)}$ carried along the ladder.
Theorem~\ref{thm:rung} and Corollary~\ref{cor:chain} are therefore statements about
the \emph{carried-solution} rung: given the ladder and its coarse solutions, the
bound is certified. A \emph{solution-free} rung --- one whose right-hand side is
computable from $(c,b,\mathrm{up})$ alone --- is not claimed. See
Proposition~\ref{prop:negative}.
\end{convention}

\begin{remark}[the convention is a scope, not a circularity]
\label{rem:scope}
It is worth being explicit about what Convention~\ref{conv:carried} does and does not
cost, because a bound whose right-hand side involves a solution invites the suspicion
that it presupposes what it proves. It does not: $u^{(l)}$ is the Galerkin solution on
the \emph{coarse} space, computed there, and \eqref{eq:8R} bounds the increment gained
by refining --- so nothing in $\widehat c_{l}$ requires knowing $\varepsilon_{l}$,
$\varepsilon_{l+1}$, or the fine solution $u^{(l+1)}$. This is precisely the
information an a posteriori estimator has available: hierarchical estimators in the
sense of \cite{BankSmith1993,BornemannErdmannKornhuber1996,DorflerNochetto2002}
evaluate a residual against the computed coarse solution and are two-sided only under
a saturation hypothesis, whereas Theorem~\ref{thm:rung} supplies the lower direction
unconditionally, at the price of being a statement \emph{given} the ladder. What the
convention rules out is a different and stronger object --- an a priori bound in the
data alone --- and Proposition~\ref{prop:negative} measures how far that object is
from reach.
\end{remark}

\section{The composed chain}
\label{sec:chain}

Theorem~\ref{thm:rung} bounds one rung; this section is its application. Two further
ingredients turn the rungs into a statement that travels: the telescope of
Lemma~\ref{lem:telescope}, which needs only semidefiniteness at the finest level, and
a bordering step that grows the window. The composition below is conditional, and on
the instance of Section~\ref{sec:numerics} its hypotheses are certified window by
window rather than proved for all windows; item~\ref{L:uniform} of
Section~\ref{sec:limits} says what would be needed for the latter.

\subsection{Windows at one common resolution}

\begin{definition}[admissible common resolution]
\label{def:admissible}
Let $u_{i}<u_{i+1}<\dots<u_{j}$ be consecutive comb nodes with log-gaps
$g_{k}=u_{k+1}-u_{k}$, and fix $\nu\ge4$. Put
\begin{equation}
\label{eq:D}
  D=\frac{\min_{k}g_{k}}{2\nu},
  \qquad
  M_{k}=\min\bigl\{M\in2\N:\;MD>u_{k}+D\bigr\},
  \qquad
  \alpha_{k}=\frac{M_{k}D}{2},
\end{equation}
and let $A^{(k)},b^{(k)},\varepsilon_{k},Q^{(k)}$ be the objects of
Definition~\ref{def:window} at $(D,M_{k})$ with the lag sequence \eqref{eq:lags}.
We call $D$ \emph{admissible} for the run and write $g_{c}^{(k)}=(M_{k+1}-M_{k})/2$
for the window increment per end.
\end{definition}

The resolution is common to the whole run \emph{by construction}, and this is the
point: distinct nodes carry distinct natural resolutions $g_{k}/(2\nu)$, and those do
not refine onto one another (a ratio $g_{k}/g_{k'}$ is a dyadic power only
exceptionally --- on the verified instance no consecutive gap ratio is dyadic, see
Section~\ref{sec:numerics}). Choosing one $D$ for the whole run buys an exactly
nested family of windows at the price of the run length, since $D$ is fixed by the
\emph{smallest} gap it contains.

\begin{lemma}[the incoming node contributes exactly nothing]
\label{lem:zeroblock}
Let $D$ be admissible and let $u$ be a node with $u>2\alpha_{k}=M_{k}D$ (i.e.\ a node
not admitted at window $k$). Then
\begin{equation}
\label{eq:zero}
  S_{D}(mD,u)=0\qquad\text{for all }0\le m\le M_{k}-1 .
\end{equation}
Consequently $c^{(k+1)}_{m}=c^{(k)}_{m}$ for all $0\le m\le M_{k}-1$.
\end{lemma}

\begin{proof}
$S_{D}(s,u)=\tfrac12[\Delta_{D}(s-u)+\Delta_{D}(s+u)]$ is supported in
$\{|s-u|<D\}\cup\{|s+u|<D\}$. For $s=mD\ge0$ and $u>0$ we have $s+u\ge u\ge\log 2>D$
on the run (indeed $u\ge u_{i}$ and $D\le\log2/(2\nu)$), so the second hat vanishes.
For the first, $m\le M_{k}-1$ gives $s\le (M_{k}-1)D$, hence
$u-s\ge M_{k}D-(M_{k}-1)D=D$ and $\Delta_{D}(s-u)=0$. The consequence is immediate
from \eqref{eq:lags}, because the two windows admit the same nodes with
$u_{j}\le 2\alpha_{k}$, all newly admitted nodes satisfy $u>M_{k}D$, and $a(mD)$ does
not depend on $M$.
\end{proof}

\begin{lemma}[exact handover identity]
\label{lem:handover}
Let $D>0$, let $M_{o}<M_{n}$ be even with $M_{n}-M_{o}=2g_{c}$, and let
$c^{(o)},c^{(n)}$ be lag sequences with $c^{(n)}_{m}=c^{(o)}_{m}$ for
$0\le m\le M_{o}-1$. Let $A^{(o)},b^{(o)}$ and $A^{(n)},b^{(n)}$ be as in
Definition~\ref{def:window} at $(D,M_{o})$ and $(D,M_{n})$, and write $h_{o}=M_{o}/2$.
Then, with indices $r=g_{c}+r'$, $s=g_{c}+s'$ for $r',s'\in\{0,\dots,h_{o}-1\}$,
\begin{equation}
\label{eq:handover}
  A^{(n)}_{rs}=A^{(o)}_{r's'},
  \qquad
  b^{(n)}_{r}=b^{(o)}_{r'} .
\end{equation}
Hence the trailing $h_{o}\times h_{o}$ block of $Q^{(n)}=A^{(n)}-b^{(n)}(b^{(n)})^{\T}$
equals $Q^{(o)}$ exactly.
\end{lemma}

\begin{proof}
For the form, $A^{(n)}_{rs}=c^{(n)}_{|r-s|}-c^{(n)}_{M_{n}-1-r-s}$ by
\eqref{eq:Adef}. Now $|r-s|=|r'-s'|\le h_{o}-1\le M_{o}-1$ and
$M_{n}-1-r-s=M_{n}-2g_{c}-1-r'-s'=M_{o}-1-r'-s'$, which lies in
$\{1,\dots,M_{o}-1\}$ because $r'+s'\le 2h_{o}-2=M_{o}-2$. Both indices are therefore
$\le M_{o}-1$, so the hypothesis on the lag sequences applies and the right-hand side
is $c^{(o)}_{|r'-s'|}-c^{(o)}_{M_{o}-1-r'-s'}=A^{(o)}_{r's'}$. For the pole vector,
$\alpha_{n}-g_{c}D=\bigl(\tfrac{M_{n}}2-g_{c}\bigr)D=\tfrac{M_{o}}2D=\alpha_{o}$, so
by \eqref{eq:cells}
$\bar x^{(n)}_{r}=-\alpha_{n}+(g_{c}+r'+\tfrac12)D=-\alpha_{o}+(r'+\tfrac12)D
=\bar x^{(o)}_{r'}$, and the prefactor $8D^{-1/2}\sinh(D/4)$ in \eqref{eq:bdef} is
unchanged because $D$ is. The statement about $Q$ follows.
\end{proof}

\begin{remark}
The two lemmas together say that on an admissible run the chain
\emph{composes literally}: by Lemma~\ref{lem:zeroblock} the lag sequences agree on the
old range, so by Lemma~\ref{lem:handover} the trailing block of the next window's
bordered form is the current window's bordered form, not merely close to it. The
coordinate ordering matters and is the natural one: index $r=0$ is the outermost cell,
so a window that grows by $g_{c}$ cells per end acquires its new coordinates at the
\emph{leading} indices.
\end{remark}

\subsection{The margin-free bordering step}

\begin{lemma}[Albert's criterion]
\label{lem:albert}
Let $\mathcal M=\bigl[\begin{smallmatrix}A&C\\ C^{\T}&X\end{smallmatrix}\bigr]$ be
symmetric. Then $\mathcal M\PSD$ if and only if
\begin{equation}
\label{eq:albert}
  X\PSD,\qquad \ran(C^{\T})\subseteq\ran(X),\qquad A-CX^{+}C^{\T}\PSD .
\end{equation}
If $X\PD$ the middle condition is vacuous and $X^{+}=X^{-1}$.
\end{lemma}

\begin{proof}
\cite{Albert1969}; the range condition is Douglas' lemma \cite{Douglas1966}.
\end{proof}

The point of \eqref{eq:albert} for us is negative: \emph{nothing in it refers to the
size of $\lmin(X)$}. Contrast the norm-perturbation surrogate, which writes
$\mathcal M$ as a perturbation of $\diag(A,X)$ and asks
$\lmin(X)>\|C\|_{2}^{2}/\lmin(A)$ or the symmetric variant; that estimate divides by
$\lmin(X)$ and is therefore hostage to exactly the quantity a lower bound on
$\varepsilon$ is trying to produce. Section~\ref{sec:numerics} reports that on the
verified instance the surrogate is negative on every step while the exact Schur
complement of \eqref{eq:albert} is of order $10^{-1}$.

\subsection{The chain}

\begin{corollary}[margin-free chain composition]
\label{cor:chain}
Let $D$ be admissible for the consecutive nodes $u_{i},\dots,u_{j}$ in the sense of
Definition~\ref{def:admissible}. For each $k$ let $D=D_{0}>\dots>D_{L_{k}}$ be the
level ladder \eqref{eq:levels} at fixed $\alpha_{k}$, and assume
\begin{enumerate}[label=\textup{(A\arabic*)},leftmargin=3.2em]
\item\label{A:base} \textup{(base)} $\varepsilon^{(k)}_{L_{k}}\ge0$ at the finest
      level of the chain of window $k$ --- equivalently, by \eqref{eq:albert-eps},
      $Q^{(k,L_{k})}\PSD$. This is pure semidefiniteness: no rate, no margin.
\item\label{A:hyp} the hypotheses \ref{H:pd}--\ref{H:maj} of Theorem~\ref{thm:rung}
      at every rung of every chain, together with Convention~\ref{conv:carried};
\item\label{A:step} for every $k$, the $g_{c}^{(k)}\times g_{c}^{(k)}$ Schur
      complement $A^{(k+1)}_{\mathrm{new}}-C_{k}\,(Q^{(k)})^{-1}C_{k}^{\T}$ is
      positive semidefinite, where $A^{(k+1)}_{\mathrm{new}}$ and $C_{k}$ are the
      leading block and the coupling block of $Q^{(k+1)}$ in the splitting of
      Lemma~\ref{lem:handover}.
\end{enumerate}
Then for every $k\in\{i,\dots,j\}$
\begin{equation}
\label{eq:chain-eps}
  \varepsilon^{(k)}_{0}
  \;\ge\;\varepsilon^{(k)}_{L_{k}}+\sum_{l<L_{k}}\widehat c^{\,(k)}_{l}
  \;\ge\;\sum_{l<L_{k}}\widehat c^{\,(k)}_{l}\;>\;0,
  \qquad\text{hence}\qquad Q^{(k)}\PD ,
\end{equation}
and the run is traversed
\begin{equation}
\label{eq:chain}
  Q^{(i)}\PD\;\Longrightarrow\;Q^{(i+1)}\PD\;\Longrightarrow\;\dots\;
  \Longrightarrow\;Q^{(j)}\PD .
\end{equation}
\end{corollary}

\begin{proof}
Fix $k$. By \ref{A:hyp} and Theorem~\ref{thm:rung}, $\delta^{(k)}_{l}\ge
\widehat c^{\,(k)}_{l}>0$ for every rung; by Lemma~\ref{lem:telescope},
$\varepsilon^{(k)}_{0}=\varepsilon^{(k)}_{L_{k}}+\sum_{l}\delta^{(k)}_{l}$; and by
\ref{A:base}, $\varepsilon^{(k)}_{L_{k}}\ge0$. This is \eqref{eq:chain-eps}, and
$Q^{(k)}\PD$ follows from \eqref{eq:albert-eps} with strict inequality.
For \eqref{eq:chain}, Lemmas~\ref{lem:zeroblock} and \ref{lem:handover} put
$Q^{(k+1)}$ in the block shape
$\bigl[\begin{smallmatrix}A^{(k+1)}_{\mathrm{new}}&C_{k}\\
 C_{k}^{\T}&Q^{(k)}\end{smallmatrix}\bigr]$
with the trailing block \emph{equal} to $Q^{(k)}$. Apply Lemma~\ref{lem:albert} with
$X=Q^{(k)}$: the first condition is the previous step's conclusion, the second is
vacuous because $Q^{(k)}\PD$, and the third is \ref{A:step}.
\end{proof}

\begin{remark}[where the induction runs, and why it must]
\label{rem:direction2}
The telescope \eqref{eq:sat} places its certified bound on the \emph{coarse} level:
$\varepsilon_{\text{coarse}}=\varepsilon_{\text{fine}}+\delta$ with $\delta\ge0$, so
the target window must be the coarsest level of its chain and the base case
\ref{A:base} sits at the finest. A bound at the fine level would require certified
\emph{upper} bounds on the rungs, i.e.\ a strengthened Cauchy--Schwarz constant, which
is precisely the quantity the construction is built to avoid. This is a structural
feature, not a choice.
\end{remark}

\section{A measured obstruction: the solution-free rung}
\label{sec:negative}

Convention~\ref{conv:carried} is the one hypothesis of Corollary~\ref{cor:chain} that
is neither an identity nor a Cholesky certificate, so it deserves to be examined
rather than declared. The natural attempt to remove it is to replace the optimal test
vector $v^{\star}=U^{-1}\shat$, which needs $u^{(l)}$, by the fully explicit
$w=U^{-1}s$, which does not.

\begin{proposition}[the solution-free rung, and how it fails]
\label{prop:negative}
With $w=U_{l}^{-1}s_{l}$, the test vector $w$ in \eqref{eq:maximum} gives the valid
lower bound
\begin{equation}
\label{eq:cbexpl}
  \delta_{l}\;\ge\;\widehat c^{\,\mathrm{expl}}_{l}
  \;=\;\frac{(\shat_{l}^{\T}w)^{2}}{s_{l}^{\T}w}
  \;=\;(1-\rho_{l})^{2}\,\bigl(s_{l}^{\T}U_{l}^{-1}s_{l}\bigr),
  \qquad
  \rho_{l}=\frac{(u^{(l)})^{\T}B_{x}w}{s_{l}^{\T}U_{l}^{-1}s_{l}},
\end{equation}
in which the factor $s_{l}^{\T}U_{l}^{-1}s_{l}$ is computable from
$(c,b,\mathrm{up})$ alone and the coarse solution enters through the single scalar
$\rho_{l}$. Consequently the rung is solution-free up to a constant \emph{if and only
if} $\rho_{l}$ admits a bound away from $1$.

On the verified instance of Section~\ref{sec:numerics} it does not, and the failure is
by cancellation rather than by a small margin:
\begin{enumerate}[label=\textup{(\roman*)},leftmargin=2.6em]
\item $|\rho_{l}|$ is measured in $[0.3888,\,14.8018]$ over the $88$ rungs of the
  assembled ladder ($[0.88,\,2.61]$ on the separate level chain of the same family),
  so $\rho_{l}$ is not a perturbation of $0$;
\item $\|\shat_{l}\|^{2}/\|s_{l}\|^{2}$ is measured in
  $[6.2\cdot10^{-1},\,1.8\cdot10^{5}]$ --- below $1$ at small $\alpha$ and above it at
  large $\alpha$ --- so the coupling term does not act as a fraction of the explicit
  high-pass part at all;
\item the norm-only variant, which replaces $\shat_{l}$ by
  $\|s_{l}\|-\|B_{x}^{\T}u^{(l)}\|$ under Cauchy--Schwarz, returns exactly
  $0.0$ of $\delta_{l}$ on every one of the $88$ rungs.
\end{enumerate}
\end{proposition}

\begin{proof}[Proof of the algebraic part]
$w$ is a test vector in \eqref{eq:maximum} after replacing $S$ by $U_{l}$, which is
legitimate by the chain $0\prec S\preceq U_{l}$ established in the proof of
Theorem~\ref{thm:rung}; hence
$\delta_{l}\ge(\shat_{l}^{\T}w)^{2}/(w^{\T}U_{l}w)$ and
$w^{\T}U_{l}w=s_{l}^{\T}U_{l}^{-1}s_{l}$. Substituting
$\shat_{l}=s_{l}-B_{x}^{\T}u^{(l)}$ gives
$\shat_{l}^{\T}w=s_{l}^{\T}U_{l}^{-1}s_{l}-(u^{(l)})^{\T}B_{x}w
=(1-\rho_{l})\,s_{l}^{\T}U_{l}^{-1}s_{l}$, and \eqref{eq:cbexpl} follows.
\end{proof}

The items (i)--(iii) are measurements on a finite ladder and are reported as such;
they do not prove that no solution-free rung exists. What they do establish is the
\emph{shape} of the obstruction, and that shape is informative: since $\rho_{l}$ is of
order one and larger, the coupling term $B_{x}^{\T}u^{(l)}$ cancels the explicit
high-pass part $s_{l}$ rather than perturbing it. A solution-free rung therefore needs
a bound on the \emph{direction} of $\shat_{l}$, not on its size, and any estimate that
proceeds through norms is provably worthless here --- item (iii) is that statement,
measured.

\section{Numerical verification}
\label{sec:numerics}

\subsection{Floating-point convention}

Every positivity statement below is a completed Cholesky factorisation relative to a
declared floor. With $u=2^{-53}=1.110\cdot10^{-16}$ the unit roundoff and
$c_{h}=(h+1)/\bigl(1-(h+1)u\bigr)$, a completed Cholesky factorisation of
$A-\delta I$ with $\delta=c_{h}u\|A\|_{\infty}$ certifies $\lmin(A)>0$
\cite{Wilkinson1963,Higham2002}; the cheap certified bound
$\|A\|_{2}\le\|A\|_{\infty}=\max_{r}\sum_{s}|A_{rs}|$ is used for the norm
\cite{Gershgorin1931}. At the largest size occurring below, $h=1495$, the floor is
$1.66\cdot10^{-13}\,\|A\|_{\infty}$. Quantities obtained from a symmetric eigenvalue
routine are labelled \emph{measured} and are never used as certificates; every
reported \emph{head-room} is a factor over the declared floor, so any value above $1$
certifies the corresponding sign.

\subsection{The verified instance}

\begin{table}[t]
\centering
\small
\begin{tabularx}{\textwidth}{@{}lY@{}}
\toprule
quantity & value \\
\midrule
common resolution & $D=3.472446\cdot10^{-3}$, from
  $\min_{k}g_{k}=0.027780$ and $\nu=4$; effective
  $\nu_{\mathrm{eff}}=g_{k}/(2D)=4.0\dots58.4$ \\
run of nodes & $30$ consecutive prime-power nodes, $n=2,\dots,79$;
  $\alpha_{k}=0.3507\dots2.1494$; $h_{k}=101\dots619$;
  $g_{c}^{(k)}=4\dots58$ cells per end \\
gap facts used & $g_{k}\le\log2$ (max.\ observed $0.405465$) and
  $g_{k}\ge\log(1+1/n_{k})$ (max.\ observed $1/g_{k}=1.4985\cdot10^{5}$) \\
window nesting & $M_{k+1}^{\text{(old)}}=M_{k}^{\text{(new)}}$ on all $29$
  consecutive pairs, an integer identity \\
non-dyadic gaps & $0$ of $84$ consecutive gap ratios $g_{k}/g_{k+1}$ is a dyadic
  power: per-node resolutions have no common refinement \\
ladder size & $52$ windows ($30$ composed $+$ $22$ at per-node resolutions, $51$
  distinct $(n,D)$ pairs), $88$ rungs, chains of $2$--$5$ levels, plus $10$
  bordering-only windows up to $n=1331$ \\
largest factorisation & $h=1495$ (cap $1500$); finest chain level $h\le1400$;
  base-case sizes $h=416\dots1384$ \\
certificates & $430$ completed Cholesky factorisations, $440$ verified identities,
  all completing \\
runtime & $8.2$\,s \\
\bottomrule
\end{tabularx}
\caption{The verified instance. All entries are outputs of the single script named in
Section~\ref{sec:repro}.}
\label{tab:instance}
\end{table}

Table~\ref{tab:instance} lists the instance. Table~\ref{tab:results} lists what was
certified on it. We highlight five entries.

\emph{Identities.} The nesting \eqref{eq:nesting}, proved in
Lemmas~\ref{lem:nesting} and \ref{lem:classconsistent}, is reproduced to relative
$3.0\cdot10^{-14}$ for the form and $6.2\cdot10^{-16}$ for the pole vector on all
$88$ rungs; since it is an identity, these residuals compare two orders of summation
for the same exact quantity and are a check on the assembly code, not on the
statement. The saturation identity, the dual-residual identity \eqref{eq:dual} and
Galerkin orthogonality hold to $4.7\cdot10^{-7}$; the telescope
$\sum_{l}\delta_{l}=\varepsilon_{0}-\varepsilon_{L}$ to $3.3\cdot10^{-7}$; the
closed-form piecewise-constant lag formula to $2.2\cdot10^{-14}$. The bit-exact part
of Lemma~\ref{lem:zeroblock} is bit-exact in the run: the incoming node's
contribution restricted to the old lag range has maximum absolute entry
$0.0\cdot10^{0}$, and the trailing block of $Q^{(k+1)}$ agrees with $Q^{(k)}$ to
relative $2.9\cdot10^{-14}$, which is the roundoff of two different summation orders
for the same exact quantity (Lemma~\ref{lem:handover}).

\emph{The rung.} All three certificates of Theorem~\ref{thm:rung} complete on
$88$ of $88$ rungs, and $0<\widehat c_{l}\le\delta_{l}$ throughout with retention
$\widehat c_{l}/\delta_{l}\in[0.90709,\,0.99708]$. The direction control of
Remark~\ref{rem:direction} is also verified: the primal form \eqref{eq:primal},
evaluated at the same test vector, returns values \emph{above} $\delta_{l}$ on $88$
of $88$ rungs, i.e.\ it is an upper bound and can never serve as a supply.

\emph{The lower bound.} $\varepsilon^{(k)}_{0}\ge\sum_{l}\widehat c^{\,(k)}_{l}>0$ on
all $52$ windows. Two normalisations are reported, and the second is the informative
one. Against the measured $\varepsilon_{0}$ the retention is
$0.0013\dots0.9694$; but an $L$-level chain cannot supply more than it telescopes, so
the relevant ceiling is $1-\varepsilon_{L}/\varepsilon_{0}$, measured at
$0.001295\dots0.993502$, and against that ceiling the certified fraction is
$0.9140\dots0.9929$ on every window. The certified margin exceeds the \emph{declared}
floating-point bound $c_{h}u\cond(A^{(k)}_{0})$ on $\varepsilon$ by a factor
$31.6\dots8.7\cdot10^{10}$ at $\cond(A_{0})=6.9\cdot10^{1}\dots3.8\cdot10^{5}$. The
thinnest window of the ladder is $n=243$ ($\alpha=2.7492$, $h_{0}=679$) at $31.6$,
and it is thin because its retention $0.6054$ is pinned to the small ceiling
$0.629042$ there, not because the arithmetic is marginal.

\emph{The base case.} At the finest level of each chain, completed Cholesky
factorisations of $A^{(L)}$ and of $Q^{(L)}=A^{(L)}-b^{(L)}(b^{(L)})^{\T}$ both run,
at sizes $h=416\dots1384$, with $\lmin(Q^{(L)})$ at
$4.4\cdot10^{5}\dots5.6\cdot10^{10}$ times its own floor; and
$\operatorname{sign}\lmin(Q^{(L)})=\operatorname{sign}\varepsilon_{L}$ on $52$ of
$52$ windows, which is \eqref{eq:albert-eps} verified rather than assumed.

\emph{The bordering step.} On $62$ of $62$ steps ($n=2,\dots,1331$,
$h_{\mathrm{new}}=30\dots1495$) the trailing block is positive definite, the range
condition of Lemma~\ref{lem:albert} is vacuous, and the exact
$g_{c}\times g_{c}$ Schur complement has
$\lmin\in[1.4839\cdot10^{-3},\,1.2528\cdot10^{-1}]$, a factor
$3.0\cdot10^{10}\dots4.6\cdot10^{13}$ above its own floor. The norm-perturbation
surrogate of the \emph{same} Schur complement, which divides by
$\lmin$ of the trailing block, is negative on all $62$ steps
($-1.04\cdot10^{7}\dots-1.29\cdot10^{1}$). Lemma~\ref{lem:albert} never performs
that division.

\begin{table}[t]
\centering
\small
\begin{tabularx}{\textwidth}{@{}lY@{}}
\toprule
statement & certified / measured on the instance \\
\midrule
$P^{\T}A^{(l+1)}P=A^{(l)}$; $P^{\T}b^{(l+1)}=b^{(l)}$
  & identity (Lem.~\ref{lem:nesting}, \ref{lem:bnesting}, \ref{lem:classconsistent});
    reproduced to $3.0\cdot10^{-14}$ / $6.2\cdot10^{-16}$ on $88/88$ rungs \\
$\varepsilon_{0}=\varepsilon_{L}+\sum\delta_{l}$
  & identity; residual $3.3\cdot10^{-7}$ \\
$\delta_{l}\ge\widehat c_{l}>0$ (Thm.~\ref{thm:rung})
  & three completed Choleskys on $88/88$ rungs;
    $\widehat c_{l}/\delta_{l}=0.90709\dots0.99708$ \\
primal form is an upper bound (Rem.~\ref{rem:direction})
  & above $\delta_{l}$ on $88/88$ rungs \\
$\varepsilon_{0}\ge\sum\widehat c_{l}>0$
  & $52/52$ windows; $0.9140\dots0.9929$ of the $L$-level ceiling; head-room
    $31.6\dots8.7\cdot10^{10}$ over the declared floor \\
base case $Q^{(L)}\PSD$
  & Cholesky at $h=416\dots1384$; head-room
    $4.4\cdot10^{5}\dots5.6\cdot10^{10}$ \\
bordering step (Lem.~\ref{lem:albert})
  & $62/62$ steps; exact Schur complement
    $\lmin=1.48\cdot10^{-3}\dots1.25\cdot10^{-1}$, head-room
    $3.0\cdot10^{10}\dots4.6\cdot10^{13}$ \\
norm-perturbation surrogate
  & negative on $62/62$ steps \\
solution-free rung (Prop.~\ref{prop:negative})
  & $|\rho_{l}|=0.3888\dots14.8018$; norm-only bound returns $0.0$ of the rung on
    $88/88$ \\
\bottomrule
\end{tabularx}
\caption{What was certified, and what was measured. Head-room is always a factor over
the declared floating-point floor of Section~\ref{sec:numerics}.}
\label{tab:results}
\end{table}

\subsection{Two fitted trends, labelled as fits}

The following are least-squares fits with jackknife bands. They are reported for
orientation only and are \emph{not} bounds; nothing in
Theorem~\ref{thm:rung} or Corollary~\ref{cor:chain} depends on them.
\begin{align}
\label{eq:fit1}
  \log\bigl(\widehat c_{l}/\delta_{l}\bigr)
  &= -0.082-0.010\,\log D+0.001\,\log\alpha,
  && \text{rms }0.012,\ \text{bands }\pm0.004,\ \pm0.005;\\
\label{eq:fit2}
  \log\Bigl(\textstyle\sum_{l}\widehat c_{l}\big/\sum_{l}\delta_{l}\Bigr)
  &= -0.018-0.012\,\log\alpha,
  && \text{rms }0.004,\ \text{band }\pm0.002 .
\end{align}
Equation~\eqref{eq:fit2} is fitted on the composed run only, where $D$ is common by
construction so no exponent in $D$ is identifiable. Read together, they say that on
this ladder the certified fraction of the telescope does not thin out along the run.

\subsection{The near-null direction}
\label{sec:nearnull}

The measurement that motivates Section~\ref{sec:intro}'s claim about symbol methods is
the following, taken on the same family of forms at the two-level stage. Let $\ell$ be
a certified pointwise minorant of the symbol, $\ell\le\sigma^{(M)}$, so that
$\Tm{M}(\ell)\preceq\Tm{M}(c)$ by Lemma~\ref{lem:parseval}, and let
$E_{c}=V^{\T}B^{\T}\Tm{M}(\ell)BV$ be the coarse block it produces. Then
$\lmin(E_{c})<0$ on every tested row, so the minorant certifies nothing about
$A_{c}$; whereas $\lmin(A_{c})$ itself is positive, measured in
$[1.7\cdot10^{-5},\,2.9\cdot10^{-4}]$, at condition numbers up to
$3.1\cdot10^{7}$ --- two to five orders of magnitude below the size the minorant
would have to reproduce. The near-null eigenvector of $A_{c}$ is smooth and lies in
$\ran V$: its small Rayleigh quotient is produced by cancellation between the
Toeplitz entries, not by a small value of the symbol along it, and a pointwise
minorant is by construction blind to that. This is exactly the quantity that
Theorem~\ref{thm:rung} does not need, since it uses $A_{c}$ only through
$A_{c}\PSD$.

\subsection{Reproducibility}
\label{sec:repro}

All numbers in Tables~\ref{tab:instance} and \ref{tab:results}, and item~(iii) and the
first band of item~(i) of Proposition~\ref{prop:negative}, are produced by one
self-contained Python script (\texttt{grand\_assembly\_probe.py}, supplied as
supplementary material) using only \texttt{numpy} and \texttt{scipy}
(\texttt{cho\_factor}, \texttt{eigvalsh}). It performs no file writes, runs in about
$8.2$\,s on a current laptop, and reports $34$ internal checks, all passing. Every
matrix that is factorised, inverted or diagonalised has $h\le1500$; the only larger
objects are matrix-free fast Fourier transforms used to evaluate symbols, at grid
size up to $2^{20}$. The two remaining sets of numbers --- the second band of
Proposition~\ref{prop:negative}(i), measured on a single level chain rather than on
the assembled ladder, and the two-level measurements of
Section~\ref{sec:nearnull} --- come from two further scripts of the same family,
supplied alongside it; they are reported separately for exactly that reason.

\section{Limitations, and what is not claimed}
\label{sec:limits}

\begin{enumerate}[label=\textup{(L\arabic*)},leftmargin=3.2em]
\item\label{L:finite} \textbf{The run is finite, and its length is structurally
  bounded.} By Definition~\ref{def:admissible} the common resolution is
  $D=\min_{k}g_{k}/(2\nu)$, the window sizes are $M_{k}\approx u_{k}/D$, and the cost
  of a factorisation is cubic in $M_{k}$. A longer run therefore forces a smaller $D$
  and larger windows, and the run verified here --- $30$ consecutive nodes,
  $n=2,\dots,79$ --- is what fits under a matrix cap of $1500$. Nothing is claimed
  about nodes outside the tested set: of the $34\,034$ comb nodes tabulated up to
  $n=400\,000$, exactly $48$ were touched at all, and $62$ of the $85$ admissible
  bordering steps in that table were certified.
\item\label{L:uniform} \textbf{No uniformity in the window index.} This is the main
  gap, and it is a gap in \emph{uniformity}, not in size. Specifically:
  hypothesis~\ref{A:step} of Corollary~\ref{cor:chain} is verified window by window
  and we have no uniform-in-$k$ lower bound for the $g_{c}\times g_{c}$ Schur
  complement; and we have no uniform lower bound for $\widehat c_{l}$. Without both,
  \eqref{eq:chain} does not extend to an infinite family, and we make no such claim.
  Consequently Corollary~\ref{cor:chain} should be read as we have labelled it, an
  application of Theorem~\ref{thm:rung} together with a certified instance of its
  hypotheses, and not as a theorem about the family of all windows.
  Theorem~\ref{thm:rung} itself is unaffected: it is a statement about one rung, with
  hypotheses that are checkable on that rung.
\item\label{L:base} \textbf{The base case is a certificate, not a theorem.}
  Hypothesis~\ref{A:base} is a completed Cholesky factorisation of one explicit
  matrix of size $h=416\dots1384$ per window. It is verified, not proved.
\item\label{L:conv} \textbf{Convention~\ref{conv:carried} stands.} The certified rung
  is the carried-solution rung. Proposition~\ref{prop:negative} shows that removing
  the convention requires a bound on the direction of $\shat_{l}$, and that
  norm-based attempts fail by cancellation on the tested ladder; it does not show
  that no such bound exists.
\item\label{L:secondroute} \textbf{A second, independent route is deliberately
  excluded.} The increment $\delta_{l}$ also admits a lower bound through a
  corner estimate for the oscillation of the Galerkin solution --- a one-cell
  Poincar\'e step in the spirit of \cite{PayneWeinberger1960} followed by a
  Cauchy--Schwarz reduction to the corner mass --- which requires a sign
  hypothesis on the increment sequence and a summed Harnack-type ratio bound. Both
  hold on the windows we measured and neither is proved uniformly; their classical
  address is the corner regularity of inverses of Toeplitz forms with singular symbols
  \cite{Widom1974,FisherHartwig1969,Trench1964} for the ratio, and the discrete
  maximum principle for monotone matrices \cite{Ostrowski1937,Varga1962} for the sign.
  Since that route is logically independent of Theorem~\ref{thm:rung} and
  Corollary~\ref{cor:chain}, which do not use it, we omit it here and defer it to a
  separate note whose subject is a discrete Harnack inequality for the increments of
  the inverse.
\item\label{L:fits} \textbf{Fits are fits.} Equations~\eqref{eq:fit1} and
  \eqref{eq:fit2}, and every exponent quoted for orientation, are least-squares fits
  with jackknife bands, not bounds.
\item\label{L:example} \textbf{The example class is an example.}
  Theorem~\ref{thm:rung} and Corollary~\ref{cor:chain} are symbol-agnostic; the class
  of Definition~\ref{def:class} enters only through the verification, and the
  arithmetic in it enters only as the source of the comb positions and through the two
  classical gap inequalities named in Section~\ref{sec:example}(iii). No asymptotic
  property of the symbol family and no arithmetic statement is used or claimed.
\end{enumerate}

\section*{Acknowledgements}

The construction grew out of a long sequence of failed attempts, and the two that
were most instructive are recorded above: the primal form of a rung
(Remark~\ref{rem:direction}, an upper bound where a lower bound was wanted) and the
norm-perturbation surrogate of the bordering step
(Section~\ref{sec:chain}, negative on every step). Both are reported as measurements
because both are, in retrospect, the reason the working argument has the shape it has.

\begin{thebibliography}{99}\small

\bibitem{Albert1969} A.~Albert,
\emph{Conditions for positive and nonnegative definiteness in terms of pseudoinverses},
SIAM J. Appl. Math. \textbf{17} (1969), 434--440.

\bibitem{AxelssonVassilevski1989} O.~Axelsson and P.~S.~Vassilevski,
\emph{Algebraic multilevel preconditioning methods, I},
Numer. Math. \textbf{56} (1989), 157--177.

\bibitem{BankSmith1993} R.~E.~Bank and R.~K.~Smith,
\emph{A posteriori error estimates based on hierarchical bases},
SIAM J. Numer. Anal. \textbf{30} (1993), 921--935.

\bibitem{BoettcherSilbermann2006} A.~B\"ottcher and B.~Silbermann,
\emph{Analysis of Toeplitz Operators}, 2nd ed., Springer, Berlin, 2006.

\bibitem{BornemannErdmannKornhuber1996} F.~A.~Bornemann, B.~Erdmann and R.~Kornhuber,
\emph{A posteriori error estimates for elliptic problems in two and three space
dimensions}, SIAM J. Numer. Anal. \textbf{33} (1996), 1188--1204.

\bibitem{BPX1990} J.~H.~Bramble, J.~E.~Pasciak and J.~Xu,
\emph{Parallel multilevel preconditioners},
Math. Comp. \textbf{55} (1990), 1--22.

\bibitem{Brandt1977} A.~Brandt,
\emph{Multi-level adaptive solutions to boundary-value problems},
Math. Comp. \textbf{31} (1977), 333--390.

\bibitem{Cea1964} J.~C\'ea,
\emph{Approximation variationnelle des probl\`emes aux limites},
Ann. Inst. Fourier \textbf{14} (1964), 345--444.

\bibitem{Chebyshev1852} P.~L.~Chebyshev,
\emph{M\'emoire sur les nombres premiers},
J. Math. Pures Appl. \textbf{17} (1852), 366--390.

\bibitem{DorflerNochetto2002} W.~D\"orfler and R.~H.~Nochetto,
\emph{Small data oscillation implies the saturation assumption},
Numer. Math. \textbf{91} (2002), 1--12.

\bibitem{Douglas1966} R.~G.~Douglas,
\emph{On majorization, factorization, and range inclusion of operators on Hilbert
space}, Proc. Amer. Math. Soc. \textbf{17} (1966), 413--415.

\bibitem{Durbin1960} J.~Durbin,
\emph{The fitting of time series models},
Rev. Inst. Int. Statist. \textbf{28} (1960), 233--244.

\bibitem{Fejer1916} L.~Fej\'er,
\emph{\"Uber trigonometrische Polynome},
J. Reine Angew. Math. \textbf{146} (1916), 53--82.

\bibitem{FisherHartwig1969} M.~E.~Fisher and R.~E.~Hartwig,
\emph{Toeplitz determinants: some applications, theorems, and conjectures},
Adv. Chem. Phys. \textbf{15} (1969), 333--353.

\bibitem{Gershgorin1931} S.~Gershgorin,
\emph{\"Uber die Abgrenzung der Eigenwerte einer Matrix},
Izv. Akad. Nauk SSSR Ser. Mat. \textbf{6} (1931), 749--754.

\bibitem{GrenanderSzego1958} U.~Grenander and G.~Szeg\H{o},
\emph{Toeplitz Forms and Their Applications},
University of California Press, Berkeley, 1958.

\bibitem{Haynsworth1968} E.~V.~Haynsworth,
\emph{Determination of the inertia of a partitioned Hermitian matrix},
Linear Algebra Appl. \textbf{1} (1968), 73--81.

\bibitem{Higham2002} N.~J.~Higham,
\emph{Accuracy and Stability of Numerical Algorithms}, 2nd ed., SIAM,
Philadelphia, 2002.

\bibitem{KMS1953} M.~Kac, W.~L.~Murdock and G.~Szeg\H{o},
\emph{On the eigenvalues of certain Hermitian forms},
J. Rational Mech. Anal. \textbf{2} (1953), 767--800.

\bibitem{Kantorovich1948} L.~V.~Kantorovich,
\emph{Functional analysis and applied mathematics},
Uspekhi Mat. Nauk \textbf{3} (1948), 89--185.

\bibitem{Levinson1947} N.~Levinson,
\emph{The Wiener RMS error criterion in filter design and prediction},
J. Math. Phys. \textbf{25} (1947), 261--278.

\bibitem{Loewner1934} K.~L\"owner,
\emph{\"Uber monotone Matrixfunktionen},
Math. Z. \textbf{38} (1934), 177--216.

\bibitem{MontgomeryVaughan1974} H.~L.~Montgomery and R.~C.~Vaughan,
\emph{Hilbert's inequality},
J. London Math. Soc. (2) \textbf{8} (1974), 73--82.

\bibitem{Ostrowski1937} A.~M.~Ostrowski,
\emph{\"Uber die Determinanten mit \"uberwiegender Hauptdiagonale},
Comment. Math. Helv. \textbf{10} (1937), 69--96.

\bibitem{PayneWeinberger1960} L.~E.~Payne and H.~F.~Weinberger,
\emph{An optimal Poincar\'e inequality for convex domains},
Arch. Rational Mech. Anal. \textbf{5} (1960), 286--292.

\bibitem{Riesz1916} F.~Riesz,
\emph{\"Uber ein Problem des Herrn Carath\'eodory},
J. Reine Angew. Math. \textbf{146} (1916), 83--87.

\bibitem{Szego1939} G.~Szeg\H{o},
\emph{Orthogonal Polynomials},
Amer. Math. Soc. Colloquium Publications \textbf{23}, New York, 1939.

\bibitem{Trench1964} W.~F.~Trench,
\emph{An algorithm for the inversion of finite Toeplitz matrices},
J. Soc. Indust. Appl. Math. \textbf{12} (1964), 515--522.

\bibitem{Varga1962} R.~S.~Varga,
\emph{Matrix Iterative Analysis},
Prentice-Hall, Englewood Cliffs, 1962.

\bibitem{Widom1974} H.~Widom,
\emph{Asymptotic behavior of block Toeplitz matrices and determinants},
Adv. Math. \textbf{13} (1974), 284--322.

\bibitem{Wilkinson1963} J.~H.~Wilkinson,
\emph{Rounding Errors in Algebraic Processes},
Prentice-Hall, Englewood Cliffs, 1963.

\bibitem{Yserentant1986} H.~Yserentant,
\emph{On the multi-level splitting of finite element spaces},
Numer. Math. \textbf{49} (1986), 379--412.

\end{thebibliography}

\end{document}
